\chapter{Fundamentação Teórica}
\label{chap:fundamentacao}

Este capítulo apresenta os conceitos estruturais necessários para a compreensão do sistema proposto. Inicialmente, discute-se a evolução do monitoramento clássico para a observabilidade em ambientes distribuídos, detalhando os tipos de dados de telemetria gerados. Em seguida, aborda-se as limitações dos modelos tradicionais de alerta e o paradoxo da explosão de dados, justificando a necessidade arquitetural de abordagens inteligentes na resolução de incidentes.

\section{Observabilidade em Arquiteturas Nativas em Nuvem}
\label{sec:observabilidade}

Historicamente, o gerenciamento de infraestrutura de Tecnologia da Informação (TI) operava de maneira reativa, focado no monitoramento clássico de servidores físicos ou máquinas virtuais monolíticas. O objetivo principal era responder à pergunta binária: "O sistema está funcionando?". Contudo, a transição para arquiteturas nativas em nuvem (\textit{cloud-native}), compostas por centenas de microsserviços efêmeros e contêineres, tornou essa abordagem insuficiente. Nesse novo paradigma, a pergunta fundamental mudou para: "Por que o sistema não está funcionando como deveria?" \cite{wang_2025_integrating}.

A crescente complexidade desses sistemas impôs novos desafios à engenharia de operações. Nesse contexto, o conceito de observabilidade emergiu como um pilar fundamental para a compreensão do comportamento interno de sistemas a partir puramente de suas saídas externas \cite{notaro2021survey}. Diferentemente do monitoramento tradicional, baseado em painéis de controle estáticos, a observabilidade permite diagnósticos precisos e ações eficazes diante de falhas complexas.

Para que um sistema seja considerado observável, a literatura consolidada o estrutura em torno de três tipos de dados de telemetria, conhecidos como os três pilares da observabilidade: métricas, \textit{logs} e rastros (\textit{traces}) \cite{zhang2024survey}. Cada pilar captura uma dimensão distinta do sistema e, quando correlacionados, fornecem uma visão abrangente do estado da infraestrutura.

\subsection{Métricas (\textit{Metrics})}
As métricas são medições quantitativas coletadas continuamente de componentes da infraestrutura, como uso de CPU, utilização de memória, I/O de disco e taxa de requisições por segundo. Por sua natureza numérica e estruturada, constituem séries temporais que permitem a identificação de tendências, padrões sazonais e desvios, sendo extremamente eficientes em termos de armazenamento computacional.

No contexto de plataformas corporativas, as métricas representam a etapa primária de detecção de anomalias (\textit{Failure Perception}) e a principal fonte para o disparo de alertas. O Prometheus, ferramenta \textit{open source} padrão em ecossistemas de nuvem, é utilizado para a raspagem periódica de dados. Contudo, o modelo clássico adotado, baseado em violações de limiares estáticos, frequentemente falha em isolar a explicação semântica do comportamento.

\subsection{Logs}
Os \textit{logs} consistem em registros textuais discretos, imutáveis e detalhados sobre os eventos que ocorrem dentro da aplicação ou no sistema operacional, abrangendo exceções, rastreamentos de pilha (\textit{stack traces}) e tempos limite de banco de dados. Diferentemente das métricas, sua natureza é não estruturada, o que os torna a fonte mais rica de contexto real sobre a causa de uma falha.

Conforme apontam \citeonline{notaro2021survey}, a correlação dos \textit{logs} com as métricas é vital, pois uma única exceção em arquiteturas de microsserviços pode desencadear uma reação em cascata. Ferramentas como o Loki (da pilha LGTM) são responsáveis por agregar esses textos. Na prática arquitetural do sistema proposto, quando o Prometheus detecta um pico matemático anômalo, o \textit{middleware} realiza o cruzamento consultando os \textit{logs} do Loki numa janela temporal exata (ex: $\pm 15$ minutos) para extrair o motivo primário do incidente.

\subsection{Rastros (\textit{Traces})}
Em ambientes distribuídos, uma requisição transita por diversos serviços até ser totalmente processada. Os rastros mapeiam o ciclo de vida completo dessa jornada. Compostos por segmentos (\textit{spans}) que registram o tempo gasto em cada nó da rede, os rastros oferecem a visibilidade necessária para isolar em qual microsserviço exato ocorreu a latência estrutural \cite{wang_2025_integrating}.

Apesar de sua enorme importância descritiva, a integração técnica de grafos de \textit{traces} diretos para modelos LLMs ainda representa um desafio emergente na literatura de AIOps \cite{zhang2024survey}. Portanto, como foco imediato de resolução rápida, a literatura prioriza a correlação entre as métricas e os \textit{logs}.

\subsection{O Paradoxo dos Dados e a Fadiga de Alertas}
\label{subsec:limitacoes_monitoramento}

A adoção do padrão OpenTelemetry e da pilha LGTM (Loki, Grafana, Tempo, Mimir) solucionou o desafio tecnológico da coleta de telemetria \cite{wang_2025_integrating}. Entretanto, essa mesma solução fomentou um novo gargalo operacional: a explosão volumétrica de dados sem correlação nativa.

O modelo clássico lida muito mal com falhas sistêmicas dinâmicas. Um incidente crítico dispara uma alta densidade de notificações desorganizadas, causando o fenômeno comumente classificado na indústria como fadiga de alertas (\textit{alert fatigue}) \cite{notaro2021survey}. A equipe de \textit{Site Reliability Engineering} (SRE) se vê obrigada a navegar freneticamente por dezenas de painéis para decifrar a raiz do problema.

Como evidenciado pelos estudos de \citeonline{moura2024deteccao}, os algoritmos de observabilidade atuais são metodologicamente impecáveis na triagem matemática da anomalia (a detecção em tempo real), porém desprovidos de capacidade de interpretação semântica da falha (a investigação textual do erro). Como consequência, o Tempo Médio de Resolução (MTTR) explode durante a fase de Análise de Causa Raiz (RCA). É exatamente nesta lacuna investigativa que a injeção cirúrgica de Modelos de Linguagem de Grande Escala (LLMs) se posiciona como a evolução arquitetural definitiva para sistemas de operações.
